\begin{abstract}
We define a minimal language called \lang that syntactically resembles LISP and
Scheme. Functions and macros in \lang are indistinguishable; the distinction is
made only at invocation. We show that the minimal features of \lang are
sufficient to define a module system that is as expressive as existing module
systems in Scheme~\cite{r6rs}.

Well-established module systems for Scheme and Racket take extra care to handle
compile-time code execution and imports and exports of macros
\cite{module-system-for-scheme,submodules-in-racket}. In contrast, we
demonstrate that such techniques are not needed for a language such as \lang
because macros have runtime semantics and are effectively identical to
functions.

The complex metaprogramming constructs used to implement the proposed module
system typically incur a significant performance penalty when combined with
runtime semantics. We propose using partial evaluation to implicitly shift
module resolution to compile time, achieving performance comparable to that of
established module systems in other languages.
\end{abstract}
